% NAL 2339, batch 2 — Gallica views 21–40.
% Pass: Opus 5.5 (claude-opus-5-5), 2026-09-23 — first pass, unchecked against the views by a human.
%
% The catalogue gives no dating for this volume: \dating{} is left empty and no
% date is inferred. No date is written on any leaf of this batch.
%
% The views. Colour portrait scans, one page per view (about 3440–3650 ×
% 4950–5650). Leaves mounted on guards, several torn at the top; three leaves
% (v. 27–30, 37–40) were folded as letters. Long s is transcribed as s.
%
% What the twenty views hold. The catalogue's title names Fermat; within this
% batch the leaves name him once, in the heading « Mr. fermat » over v. 37, in
% a paler ink and finer stroke than the text, perhaps another hand. Nothing in
% the batch says the hand is Fermat's; nothing below attributes a hand.
%
%   v. 21–26  The end of a Latin treatise on plane and solid loci in species
%             notation, begun before this batch (v. 21 opens mid-argument on
%             the circle and ellipse): the hyperbola for « Aq + Bq ad Eq in
%             ratione data », the « difficillima » case where A in E is mixed
%             with Aq and Eq, reduced by a triangle given in species to the
%             ellipse; « Breuiter igitur & dilucidè complexi sumus quidquid de
%             Locis planis & Solidis inexplicatum veteres reliquere »; the
%             last proposition of Apollonius's plane loci, book 1; and « coro-
%             nidis loco » a locus theorem, with the example of two given
%             points N, M and the locus of I where IN² + IM² has a given ratio
%             to the triangle INM (a circle). It ends at mid-page on v. 26.
%             One posed hand, catchwords at the foot of versos (« Bq », v. 22;
%             « casus », v. 24), the writer's own marginal addition keyed by an
%             asterisk (v. 22).
%   v. 27–28  « Copie d'un escrit enuoyé par le R. Pere Marçenne a [monsieur,
%             struck] des Cartes », then « Methodus ad disquirendam maximam &
%             minimam »: A, then A + E, « Adæquentur, vt Loquitur Diophantus »,
%             divide by E, strike out E; example: divide AC at E so that the
%             rectangle AEC is greatest (B æquabitur A bis). Ends « Nec potest
%             generalior dari methodus. » The French title is the only place in
%             the batch that names a sender and a receiver, and they are
%             Mersenne and Descartes, not the author of the « escrit ».
%   v. 29–30  « De tangentibus Linearum curuarum », in the same hand, on the
%             next leaf and continuing the method (« Ad superiorem methodum »):
%             the tangent BE to the parabola Bdn, CE found double of cd. Then
%             centres of gravity, quadratures, and solids « de quibus fusè iam
%             cum Dño de Roberual egimus ». A French note in the left margin of
%             v. 29, keyed by an asterisk, speaks of the author in the third
%             person (« Il a icy nommé B ce qu'il nomme d par apres »). In grey
%             pencil, written up the left margin of v. 30, another hand: « Mr des
%             Cartes f. 347 ».
%   v. 31–35  « Appendix ad Isagogen topicam continens Solutionem problematum
%             Solidorum per locos », in a second, more calligraphic hand with
%             its own corrections: cubic and quartic equations solved by two
%             « æqualitates locales » (parabola with hyperbola, parabola with
%             circle); the two mean proportionals (« mesolabium ») by parabola
%             and hyperbola, then by two parabolas, with a reference to
%             Eutocius on Archimedes; Viète's « parapleroses » dismissed
%             (« Abeant igitur »), Viète's « cap. i. de emend. » cited; the
%             construction of all cubic and quartic problems by parabola and
%             circle. Ends on v. 35: « Qui hæc aduerterit frustrà quæstionem
%             mesolabij, trisectionis angularis, et Similes tentabit deducere ex
%             planis ». The Z's of this piece are in a blacker ink, as if
%             filled in afterwards.
%   v. 36     Blank verso (the back of v. 35, showing through). Skipped.
%   v. 37–38  Headed « Mr. fermat »: three lines given in position forming a
%             triangle AMB; the point O from which lines drawn at given angles
%             make the rectangle EOD in a given ratio to the square of OI lies
%             on a conic section; construction and proof, citing « sexta
%             decima propositione 3i Apoll. », with « Superflua præsertim tibi
%             non addimus ». Ends on v. 38 « … per demonstrationem ducentem ad
%             impossibile regressus. » Written across the lower right of v. 38,
%             an address: « Pour Mons.r Carcaui [doubtful]. rue michel le Conte
%             au milieu ». Probably the hand of v. 31–35; not settled.
%   v. 39     Blank recto (numbers only, see below; a faint figure from v. 40
%             shows through). Skipped.
%   v. 40     Its verso: the figure of v. 37–38 (triangle a m b, point o, the
%             conic through o, p, u, b), with no text. Given as a note.
%
% Continuity. v. 21 continues a text begun in batch 1. The batch ends on a
% complete piece: v. 40 holds only the figure of v. 37–38, and nothing on it
% runs on to v. 41 (not seen).
%
% Numbers on the leaves. No \folio{} is written anywhere; the coordinator
% should decide which series is the library's.
%   - Red crayon, top right of every recto, blanks included, one per leaf and
%     continuous across the three hands: 7 (v. 21), 8 (23), 9 (25), 10 (27),
%     11 (29), 12 (31), 13 (33), 14 (35), 15 (37), 16 (39, blank). It is the
%     best candidate for the library's foliation; no certificate is in this
%     batch to confirm it, so no folio is written.
%   - Grey pencil, beside or under the red, mostly one more than it at first:
%     8 (v. 21), 9 (23), 10 (25), a single stroke at v. 27, then 11 (29), 12
%     (31), 13 (33), a doubtful « 16 » or « 13 » (35), 14 (37), 15 (39). A
%     second count, not continuous as read.
%   - Ink, the first treatise only: 13 (v. 21), 15 (23), 17 (25) at top right
%     of rectos, each struck through in grey pencil, and 14 (v. 22), 16 (24)
%     at top left of versos; v. 26 has a looped figure at top left, expected
%     18, not read with certainty. A pagination by pages, the writer's or
%     copyist's, of the piece that ends on v. 26.
%   - Ink, foot of the rectos of the Appendix only: 43 (v. 31), 44 (33), 45
%     (35). Another arrangement's leaf count; not a folio.
%   - Stamps and marks: « ACQ. 8070 (LIBRI) » and « R.c. 8070 (88) » on v. 27
%     (the head of a new piece, a Libri lot number); the round Bibliothèque
%     nationale stamp on v. 27, 29, 30, 31, 35, 37, 38; « Mr des Cartes f. 347 »
%     in pencil on v. 30 (a reference, unexplained on the leaf).
%
% Notation. The species notation stays as written: vowels A, E for the
% unknowns, consonants for the given magnitudes; « Aq », « Eq », « Aqq »,
% « A cubus »; « Zpl. », « B Sol. », « Dppl. » (plano-plano) for the
% dimensions; « in » for the product; « æquatur », « æquabitur », « æquale »
% for equality, and « adæquatur », « adæquabitur » for Diophantus's
% « adæqualitas » (v. 27–30). Two small strokes after a term (« A in E'' »,
% « Bq'' ») stand for « bis » and are set as $''$. A bar over a letter or over
% a pair of letters (a line or a point of the figure, or the unknown E) is set
% with \overline where the leaf has it and is clear. The copyist of v. 31–38
% writes a small u for the point V and a small n for N; they are given as V
% and N and flagged. Abbreviations with a tilde (« oĩa », « quæõnem », « nrã »,
% « Dño ») are kept as written.
%
% hand.md: no Fermat section yet; see the closing report.

% Coordinator's reconciliation (2026-09-23), after all three batches:
%   The red crayon series, top right of every recto, blanks included, runs
%   1–6 (v. 9–19), 7–16 (v. 21–39), 17–24 (v. 41–55) without a gap, and ends
%   on the two blank leaves 23 and 24, exactly as the certificate of v. 7
%   has it (« Volume de 24 Feuillets / Les Feuillets 23 & 24 sont blancs »).
%   It is the library's foliation, and not ink: \folio{} is written from it
%   on every transcribed recto, 1r–22r. The grey pencil series beside it is
%   one higher on v. 9–19, irregular on v. 21–39, and agrees from v. 41; it is
%   recorded in notes and not used. Batches 1 and 2 first wrote no \folio{};
%   the coordinator added them. The ink pagination of the loci treatise
%   (1–18) and the foot numbers 43–45 of the Appendix are the copies' own.
%   v. 61–66 (batch 4) are binding, so batch 4 has no file.
%   Fermat is named on the leaves in three headings only: « A Domino de
%   Fermat ad Dominum De Carcaui » (v. 41), « Lettre de Mons.r De fermes A
%   Monsieur De Carcaui » (v. 49), and « Mr. fermat » over v. 37 in a paler,
%   perhaps later hand; the certificate's « Mémoires de Fermat » (v. 7) is
%   the library's. No hand is attributed.

\documentclass[11pt,a4paper]{article}
% The preamble shared by every transcription of the mathematicians' manuscripts in Gallica.
%
% It rests on one idea: **uncertainty compounds, it does not resolve.** A
% transcription that smooths over a doubtful word has destroyed the only thing
% it was made for — a reader must be able to tell, at every point, what was on
% the page and what was supplied. The apparatus macros below are therefore the
% critical apparatus, not decoration.
%
% The same commands are recognised by `scripts/render.mjs`, which produces the
% site's reading view, and by `scripts/tei.mjs`. A macro added here must be
% added there too, or rendering fails loudly — which is the wanted behaviour.
%
% Comments here are English, like the rest of the repository. The documents
% this preamble sets are French, and that is deliberate: see the skills.

\ProvidesPackage{germain}[2026/09/15 Transcriptions of mathematicians' manuscripts in Gallica]

\usepackage{amsmath,amssymb,amsthm}
\usepackage{mathrsfs}
% Kept from the parent preamble so the renderer's diagram support stays
% available; these manuscripts draw no commutative diagrams, and nothing here uses it.
\usepackage{tikz-cd}
\usepackage[english,french]{babel}
\usepackage{fontspec}

% --- Greek ------------------------------------------------------------------
% The text face stays fontspec's default, Latin Modern, which has no Greek:
% XeTeX drops every glyph it lacks with a line in the log and nothing on the
% page, and the Greek words the manuscripts quote vanished from the PDFs. So
% the Greek and Greek Extended blocks get a XeTeX character class of their own,
% and entering or leaving a run of it switches the family to CMU Serif — the
% same Computer Modern design, with polytonic Greek — and back. Only the
% family changes: a Greek word in \textit or \textbf keeps its shape and series.
% The transitions to and from every other class carry the tokens that class
% has towards an ordinary letter, so French spacing before « ; » or inside
% « » still holds after a Greek word. No grouping, so no brace or \endgroup
% can come out unbalanced; maths is left alone, since XeTeX inserts nothing
% there. Kept identical in `exercices.sty`.
\makeatletter
\newfontfamily\germainGreekfont{cmun}[
  Extension=.otf, NFSSFamily=germaingreek,
  UprightFont=*rm, ItalicFont=*ti, SlantedFont=*sl,
  BoldFont=*bx, BoldItalicFont=*bi, BoldSlantedFont=*bl]
\newXeTeXintercharclass\germain@greek
\def\germain@greekrange#1#2{%
  \@tempcnta=#1\relax
  \loop
    \XeTeXcharclass\@tempcnta=\germain@greek
    \ifnum\@tempcnta<#2\advance\@tempcnta\@ne\repeat}
\germain@greekrange{"0370}{"03FF}
\germain@greekrange{"1F00}{"1FFF}
\def\germain@greekprev{\rmdefault}
\def\germain@greekin{%
  \edef\germain@greekprev{\f@family}\fontfamily{germaingreek}\selectfont}
\def\germain@greekout{\fontfamily{\germain@greekprev}\selectfont}
\def\germain@greekjoin{%
  \XeTeXinterchartokenstate=\@ne
  \@tempcnta=\z@
  \loop
    \ifnum\@tempcnta=\germain@greek\else
      \XeTeXinterchartoks\germain@greek\@tempcnta=\expandafter{%
        \expandafter\germain@greekout\the\XeTeXinterchartoks\z@\@tempcnta}%
      \XeTeXinterchartoks\@tempcnta\germain@greek=\expandafter{%
        \the\XeTeXinterchartoks\@tempcnta\z@\germain@greekin}%
    \fi
    \ifnum\@tempcnta<4095\advance\@tempcnta\@ne\repeat}
% babel-french sets its own classes, and saves and restores the interchar
% state, each time a language is selected: join again after it does.
\addto\extrasfrench{\germain@greekjoin}
\addto\extrasenglish{\germain@greekjoin}
\AtBeginDocument{\germain@greekjoin}
\makeatother

\usepackage{geometry}
\geometry{a4paper,margin=25mm,marginparwidth=28mm}
\usepackage{marginnote}
\usepackage{xcolor}
\usepackage[normalem]{ulem}
\setlength{\emergencystretch}{3em}
\usepackage{hyperref}
\hypersetup{colorlinks=true,linkcolor=black,urlcolor=black,pdfborder={0 0 0}}

% --- Batch metadata ---------------------------------------------------------
% Taken as they stand by the rendering script, which shows them at the head of
% the reading view. They fix what the file claims to cover. The macro names are
% the parent project's, so that the scripts stay shared; \folder here means the
% volume's slug — `fr-9115` — and \pages counts Gallica views.

\newcommand{\folder}[1]{\gdef\grothFolder{#1}}
\newcommand{\batch}[1]{\gdef\grothBatch{#1}}
\newcommand{\pages}[2]{\gdef\grothFirst{#1}\gdef\grothLast{#2}}
\newcommand{\foldertitle}[1]{\gdef\grothFoldertitle{#1}}

% The holder's own shelfmark, printed as they write it: « Français 9115 ».
\newcommand{\shelfmark}[1]{\gdef\grothShelfmark{#1}}

% The dating, copied verbatim from the holder's catalogue — never inferred
% here. « XIXe siècle » is the BnF's; a pass that dates a leaf from its content
% says so in a \note{}, as its own inference.
\newcommand{\dating}[1]{\gdef\grothDating{#1}}

% The status notice, printed in the title block and repeated as a diagonal
% watermark on every page of the PDF. The manuscripts are in the public
% domain; what this line declares is not a right but a quality — an
% unverified first pass by a machine. The skills write the line; a file
% without it gets no watermark, so leaving it out is visible in the output.
\newcommand{\watermark}[1]{\gdef\grothWatermark{#1}}

\gdef\grothFolder{?}\gdef\grothBatch{}\gdef\grothFirst{?}\gdef\grothLast{?}%
\gdef\grothFoldertitle{}\gdef\grothShelfmark{}\gdef\grothDating{}\gdef\grothWatermark{}

% --- The résumé -------------------------------------------------------------
\newenvironment{resume}
  {\small\setlength{\parskip}{0.45em}}
  {\par\medskip\hrule\medskip}

% --- Keywords ---------------------------------------------------------------
% In English, closing the modernised reading's résumé: the modern vocabulary
% under which to search for what the leaves construct. The manifest extracts
% them to tag the volume — this line is the single source of the tags.
\newcommand{\keywords}[1]{\par\smallskip\noindent
  {\footnotesize\textsc{Keywords} --- \textit{#1}}\par}

% --- The critical apparatus -------------------------------------------------

% The Gallica view number, in the margin. This is the anchor that lets a
% disagreement over a formula be settled by looking at *one* image rather than
% a volume of 750. The site uses it to turn the facsimile as the transcription
% is scrolled. A view is one scanned image: usually one side of a leaf.
\newcommand{\page}[1]{%
  \marginnote{\footnotesize\color{gray!70}\textsf{v.\,#1}}%
  \label{page:#1}\ignorespaces}

% The BnF's pencilled foliation, where the leaf carries it and a pass can read
% it: « 198r ». This is what the literature cites — « ff. 198r–208v » — and
% the one line that turns a citation into a view. Recorded, never inferred:
% a folio a pass cannot read is not written.
\newcommand{\folio}[1]{%
  \marginnote{\footnotesize\color{gray!50}\textsf{f.\,#1}}\ignorespaces}

% The modernised reading's coarser counterpart to \page{}: which views a
% section is drawn from.
\newcommand{\pagerange}[2]{%
  \marginnote{\footnotesize\color{gray!70}\textsf{v.\,#1--#2}}%
  \ignorespaces}

% Illegible: never guessed.
\newcommand{\ill}{\textcolor{red!60!black}{[\,\ldots\,]}}

% A doubtful reading: the word is given, and flagged as uncertain.
\newcommand{\uncertain}[1]{\uline{#1}}

% An editorial addition — a missing letter, an implied word.
\newcommand{\add}[1]{\textcolor{blue!50!black}{[#1]}}

% Struck out by the writer. What was crossed out is part of the document.
\newcommand{\struck}[1]{\sout{#1}}

% The transcriber's note, on the state of the page or the mathematics.
\newcommand{\note}[1]{\par\smallskip{\small\color{gray!60!black}\textit{[#1]}}\par\smallskip}

% A marginal note of hers, distinct from the body of the text.
\newcommand{\marginal}[1]{\par\smallskip\noindent\hspace{1em}%
  {\small\color{gray!60!black}#1}\par\smallskip}

% --- Title ------------------------------------------------------------------
% The title block is French, like the documents themselves.

\AddToHook{shipout/background}{%
  \ifx\grothWatermark\empty\else
    \begin{tikzpicture}[remember picture, overlay]
      \node[rotate=52, scale=3.4, align=center, text=gray!18,
            font=\sffamily\bfseries]
        at (current page.center) {\grothWatermark};
    \end{tikzpicture}%
  \fi}

\AtBeginDocument{%
  \begin{center}
    {\large\bfseries Manuscrits de mathématiciens (Gallica) ---
      \ifx\grothShelfmark\empty\grothFolder\else\grothShelfmark\fi}\\[2pt]
    {\itshape\grothFoldertitle}\\[4pt]
    {\small vues \grothFirst--\grothLast
      \ifx\grothBatch\empty\else\ (lot \grothBatch)\fi}\\[2pt]
    \ifx\grothDating\empty\else
      {\small\color{gray!60!black}Datation du catalogue : \grothDating}\\[2pt]
    \fi
    \ifx\grothWatermark\empty\else
      {\small\bfseries\color{red!45!gray}\grothWatermark}\\[2pt]
    \fi
    {\footnotesize\color{gray!60!black}Transcription établie d'après les images IIIF de Gallica.
     Source gallica.bnf.fr / Bibliothèque nationale de France.\\
     Les lectures incertaines sont soulignées, les illisibles notées [\,\ldots\,],
     les ajouts entre crochets ; \textsf{v.} numérote les vues de Gallica,
     \textsf{f.} la foliotation au crayon de la BnF.}
  \end{center}
  \vspace{1em}}

\endinput


\folder{nal-2339}
\shelfmark{NAL 2339}
\batch{2}
\pages{21}{40}
\dating{}
\watermark{Lecture automatique\\non vérifiée}
\foldertitle{Fermat. Mémoires de mathématiques.}

\begin{document}

\page{21}\folio{7r}
\note{En haut à droite : « 7 » au crayon rouge ; « 13 » à l'encre, barré d'un trait de crayon gris ; au-dessous « 8 » au crayon gris (voir l'en-tête). La page continue un texte commencé avant ce lot.}

Locus erit ad circulum, vt iam diximus, licet igitur coefficientes sint eadem, modo angulus non sit rectus Locus erit ad Ellipsim.

Et licet immisceantur æquationibus homogenea sub datis \& A vel $\overline{E}$ fiet reductio eo quod iam \uncertain{vsurpauimus} artificio.

Si Aq $+$ Bq est ad Eq in ratione datâ punctum I est ad \uncertain{Hyperbolen}.
\note{« Aq $+$ Bq » : le signe est bien un « $+$ » sur le feuillet. Le dernier mot finit sur une boucle qui peut être une abréviation de la finale.}

\note{Figure : une droite verticale portant de haut en bas O, R, N, M ; en O et en N deux horizontales, O---I en haut, N---Z en bas ; une courbe part de R et monte jusqu'en I ; une verticale descend de I vers la ligne NZ.}

fiat $\overline{NO}$ parallela $\overline{ZI}$. Data ratio sit eadem quæ Bq ad quadratum NR dabitur ergo punctum R circa diametrum $\overline{RO}$, per verticem R, centrum N describatur Hyperbole, cuius applicatæ sint parallelæ NZ. Et rectangulum sub tota diametro Et $\overline{RO}$ vnâ cum $\overline{RO}$ quadrato ad quadratum $\overline{OI}$ sit in data ratione NR quadrati ad Bq. Ergo componendo Rectangulum sub MOR (posita MN æquali NR) \uncertain{m} vnâ cum quadrato $\overline{NR}$ erit ad quadratum $\overline{OI}$
\note{Après la parenthèse, en bout de ligne, un petit « m » surélevé dont le sens n'est pas vu. La phrase continue en tête de la vue 22.}

\page{22}
\note{En haut à gauche, « 14 » à l'encre (pagination, voir l'en-tête). La page continue sans rupture la phrase de la vue 21.}

Vna cum Bq in ratione data NR quadrati ad Bq. Sed rectangulum MOR vnâ cum NR quad.\textsuperscript{to} æquatur $\overline{N}$Oq siue ZI quadrato siue $\overline{E}$q. Et quadratum $\overline{OI}$ Vna cum Bq.* Ergo est vt Eq. ad Bq $+$ Aq ita NR quadratum ad Bq. Et conuertendo Bq $+$ Aq est ad Eq in ratione datâ Punctum igitur $\overline{Z}$ est ad Hyperbolen positione datam.
\marginal{* æquatur quadrato NZ siue Aq vnâ cum Bq.}
\note{L'astérisque du texte, après « Bq », renvoie à la note marginale de gauche, de la même main, séparée du texte par un trait vertical. « Punctum igitur Z » : la lettre est bien le Z surligné du texte ; l'énoncé de la vue 21 disait « punctum I ».}

Eodem quo iam vsi sumus artificio ad hanc æqualitatem reducentur oẽs quæ ab Aq Et Eq adficiuntur vnâ cum datis siue simpliciter siue misceantur ipsis homogenea sub A vel E in datas, modo Aq habeat eandem ex altera parte adfectionis notam quam Eq. Nam si sint diuersa, propositio concluditur per circulos vel Ellipses.
\note{« oẽs » : « omnes » abrégé, un tilde sur la ligne.}

Difficillima omnium æqualitatum est, quando ita misceatur Aq Et Eq. vt nihilominus homogenea quædam ab A in E adficiantur vna cum datis \&c.

Bq $-$ Aq bis æquatur A in E$''$ $+$ Eq.
\note{Les deux traits après le second E sont ceux du feuillet ; ils valent sans doute « bis » (« A in E bis »), comme « bis » après « Aq » dans le même membre.}

Addatur vtrinq; Aq vt $\overline{A}$ $+$ $\overline{E}$ sit latus alterius ex homogeneis, Ergo
\note{Au pied, à droite, la réclame « Bq », qui ouvre la vue 23.}

\page{23}\folio{8r}
\note{En haut à droite : « 8 » au crayon rouge ; « 15 » à l'encre, barré d'un trait de crayon gris ; au-dessous « 9 » au crayon gris (voir l'en-tête). La page ouvre sur la réclame « Bq » de la vue 22.}

Bq $-$ Aq æquabitur Aq $+$ Eq $+$ A in E$''$

\note{Figure : une horizontale N---Z---M ; de Z monte une verticale jusqu'en I, portant V entre Z et I ; un arc de cercle va de I à M ; de M descend une verticale jusqu'en R ; la droite NR coupe en O le prolongement de IZ sous Z. La lettre V est écrite comme un petit u.}

Pro A $+$ E Sumatur $\overline{E}$ si placet. Et ex præcedentibus circulus MI præstet propositum, hoc est MN q. (siue Bq) $-$ NZ quadrato (siue Aq) æquatur quadrato ZI, siue (quadrato ab A $+$ E) fiat $\overline{VI}$ æqualis NZ siue (A, ergo $\overline{ZV}$ æquatur $\overline{E}$. In hac autem quæstione punctum $\overline{V}$ siue extremum rectæ $\overline{E}$ tantum inquirimus. Videndum ergo Et demonstrandum ad quam lineam sit punctum $\overline{V}$. fiat MR parallela ZI, Et æqualis MN, \& iungatur NR, ad quam producta IZ incidat ad punctum $\overline{O}$. Cum NM æquatur MR ergo NZ æquabitur $\overline{ZO}$. Sed $\overline{NZ}$ æquatur VI, ergo tota VO toti ZI est æqualis. Ideoq; quadratum MN $-$ quadrato $\overline{NZ}$ æquatur quadrato $\overline{VO}$.
\note{La parenthèse ouverte avant « A, ergo » n'est pas refermée sur le feuillet. Le texte s'arrête au milieu de la page ; la suite est en tête de la vue 24.}

\page{24}
\note{En haut à gauche, « 16 » à l'encre (pagination, voir l'en-tête). La page continue la démonstration de la vue 23 ; la figure de la vue 23 se voit par transparence.}

Datur autem triangulum NMR specie. Ergo quadrati $\overline{N}$M ad quadratum NR datur ratio. Ideoq; et quadrati VZ ad quadratum NO dabitur ratio. Ratio igitur quadrati MN $-$ quadrato NZ ad quadratum VR $-$ quadrato NO datur. Probauimus autem quadratum OV æquari quadrato MN $-$ quadrato NZ. Ergo ratio NR quadrati $-$ NO quadrato ad quadratum OV datur. Dantur autem puncta N et R. Et angulus NOZ. Ergo punctum $\overline{V}$ ex superius demonstratis est ad Ellipsim.
\note{Au-dessus de « VZ » et de « VR » (lignes 3 et 5), une encre plus noire et un trait plus appuyé récrivent « NZ » et « NR », sans barrer les leçons de la ligne : on les donne ici telles qu'elles sont dans la ligne. Le sens demande N, et le V (petit u) peut n'être qu'un N mal formé que le correcteur a précisé.}

Non dissimili methodo ad superiores casus reducentur reliqui in quibus homogenea sub A in E homogeneis partim datis, partim sub Aq aut Eq immiscebuntur aut etiam sub A Et $\overline{E}$ in datas ductis, cuius rei disquisitio facillima. Semper enim beneficio trianguli specie noti construitur quæstio.

Breuiter igitur \& dilucidè complexi sumus quidquid de Locis planis \& Solidis inexplicatum veteres reliquere.

Constabitq; deinceps ad quem Locum pertinebunt
\note{Au pied, à droite, la réclame « casus », premier mot de la vue 25, où la phrase continue.}

\page{25}\folio{9r}
\note{En haut à droite : « 9 » au crayon rouge ; « 17 » à l'encre, barré d'un trait de crayon gris ; au-dessous « 10 » au crayon gris (voir l'en-tête). La page ouvre sur la réclame « casus » de la vue 24.}

casus omnis propositionis vltimæ lib. 1 Apoll. de Locis planis, Et oĩa oĩo ad hanc materiam spectantia nullo negotio detegentur.
\note{« oĩa oĩo » : « omnia omnino » abrégés.}

Sed libet coronidis loco pulcherrimam hanc propositionem adiungere, cuius facilitas statim innotescet.

Si positione datis quotcunq; lineis, ab vno \& eodem puncto. ad singulas ducantur rectæ in datis angulis, Et sint species ab oĩbus ductis dato spatio æqualis. Punctum continget positione datũ solidum Locum. Vnico exemplo fit via ad practicam.

\note{Figure : une horizontale portant N, A, Z, M ; de Z s'élève une perpendiculaire, marquée E vers son pied, jusqu'à un point en haut, dont la lettre n'est qu'un trait (le I du texte). Un grand cercle très pâle, au crayon, passe autour de la figure ; il n'est pas de l'encre du texte.}

Datis duobus punctis NM, inueniendus Locus a quo si iungas rectas IN, IM. Quadrata rectarum IN, IM ad triangulum INM habeant datam rationem. Recta NM æquetur B. Et recta ZI ad angulos rectos dicatur $\overline{E}$, terminus NZ dicatur $\overline{A}$. Ergo ex artis præceptis Aq bis $+$ Bq $-$ B in A$''$ $+$ Eq bis ad rectangulum B in E habebit rationem datam. Et resoluendo hypostasis ex iam traditis præceptis ita procedet constructio.
\note{Le texte s'arrête aux deux tiers de la page, sur ces mots ; la construction annoncée est à la vue 26.}

\page{26}
\note{En haut à gauche, un nombre à l'encre d'un seul trait bouclé, qu'on attend « 18 » dans la pagination (voir l'en-tête) mais qui ne se lit pas sûrement : on ne le donne pas.}

\note{Figure : un demi-cercle (ou un arc) passant par V en haut, O à droite et R à gauche, au-dessus d'une horizontale N---Z---M ; de Z monte une verticale, marquée I là où elle coupe l'arc inférieur, jusqu'en V ; V est joint à O ; R est joint à N, et deux droites se croisent entre Z et O, l'une allant de Z vers O, l'autre de R vers M.}

NM bifariam secetur in Z, a puncto Z. excitetur perpendicularis ZV. Et fiat data ratio eadem quæ ZV quadrupla ad ZM. descripto semicirculo VOZ super $\overline{VZ}$ applicetur \struck{ipsi} ZO æqualis ipsi \struck{ZO} ZM, Et iunctâ VO, Centro $\overline{V}$, interuallo VO describatur circulus OIR. In quo sumatur quodlibet punctum vt $\overline{R}$ Et iungantur rectæ RN, RM: Aio quadrata RN, RM ad triangulum RNM esse in data ratione.
\note{La correction de la ligne 4 : « ipsi » est barré devant « ZO » ; au-dessus de « ZO » est écrit puis barré « ZM » ; au-dessus de « æqualis ZM » est ajouté « ipsi » suivi d'un « ZO » barré. La leçon finale est « applicetur ZO æqualis ipsi ZM ». Le texte s'arrête au milieu de la page, et le reste est blanc : c'est la fin de la pièce commencée avant ce lot.}

\page{27}\folio{10r}
\note{Nouvelle pièce, sur un nouveau feuillet. En haut à droite, « 10 » au crayon rouge, et à côté un seul trait oblique au crayon gris, qui peut être un « 1 » ; aucun chiffre à l'encre à droite. En haut à gauche, un signe à l'encre en forme de croix à bras recourbés (un « Y » ou une croix), qui n'est pas un nombre. Cachets « ACQ. 8070 (LIBRI) » en tête et « BIBLIOTHEQUE NATIONALE R.F. MANUSCRITS » dans la marge gauche ; au pied, à l'encre d'une autre main, « R.c. 8070 (88) ».}

Copie d'un escrit enuoyé par le R. Pere Marçenne a \struck{\uncertain{monsieur}} des Cartes
\note{Ce titre en français, sur deux lignes au-dessus du titre latin, est de la même encre brune que le texte ; le C initial est une grande capitale ornée. C'est le seul endroit du lot qui nomme quelqu'un en titre : il nomme Mersenne et Descartes, non l'auteur de l'écrit.}

\section*{Methodus ad disquirendam maximam \& minimam.}

Omnis de inuentione maximæ \& minimæ doctrina duabus positionibus ignotis innititur, \& hac vnicâ præceptione.

Statuatur quilibet quæstionis terminus esse A siue planum siue solidum aut Longitudo, prout proposito satisfieri par est, Et inuentâ maximâ aut minimâ in terminis sub A gradu vtlibet inuolutis. Ponatur rursus idem qui prius terminus esse A $+$ E, iterumq; inueniatur maxima aut minima in terminis sub A \& E gradibus vtlibet coefficientibus. Adæquentur, vt Loquitur Diophantus, duo homogenea maximæ aut minimæ æqualia Et demptis communibus, quo peracto homogenea omnia ex parte alterutrâ (ab E vel ipsius gradibus afficiuntur) applicentur omnia ad E vel ad elatiorē ipsius gradum, donec aliquod ex homogeneis ex parte vtrauis adfectione sub E omnino liberetur. Elidantur deinde vtrinq; homogenea sub E aut ipsius gradibus quomodolibet inuoluta \& reliqua æquentur. Aut si ex vna parte nihil superest,
\note{La phrase continue en tête de la vue 28.}

\page{28}
\note{Verso du feuillet de la vue 27 ; aucun nombre. Le cachet et la cote de la vue 27 se voient par transparence.}

Adæquentur sane, quod eodem recidit negata affirmatis. Resolutio ultima istius æqualitatis dabit valorem A, quâ cognitâ, maxima aut minima ex repetitis prioris resolutionis vestigiis innotescet.

Exemplum subiicimus.

\note{Figure : un segment A---E---C, E entre A et C, plus près de A.}

Sit recta AC ita diuidenda in $\overline{E}$ vt rectangulum $\overline{AEC}$ sit maximum.

Recta $\overline{AC}$ dicatur $\overline{B}$, ponatur pars altera B esse $\overline{A}$, ergo reliqua erit B $-$ A. Et rectangulum sub segmentis erit B in A $-$ Aq. Quod debet inueniri maximum. Ponatur rursus pars altera ipsius B esse A $+$ E, ergo reliqua erit B $-$ A $-$ E. Et rectangulum sub segmentis erit B in A $-$ Aq $+$ B in E $-$ A in E bis $-$ Eq. quod debet adæquari superiori rectangulo B in A $-$ Aq. demptis communibus, B in E adæquabitur A in E bis $+$ Eq et omnibus per E diuisis B adæquabitur A bis $+$ E. elidatur E. B æquabitur A bis. Igitur B bifariam est diuidenda ad solutionem propositi
\note{« B in A $-$ Aq » (ligne 6 du paragraphe) : un petit rond est tracé au-dessus du signe moins ; il n'en change pas la lecture, que les lignes suivantes répètent avec un moins simple.}

Nec potest generalior dari methodus.
\note{Deux traits obliques suivent le mot « methodus » ; le reste de la page est blanc. C'est la fin de la pièce ouverte à la vue 27.}

\page{29}\folio{11r}
\note{En haut à droite : « 11 » au crayon rouge ; au-dessous « 11 » au crayon gris, en deux traits obliques (voir l'en-tête). Aucun nombre à l'encre. Cachet rond de la Bibliothèque nationale dans la marge gauche, au pied.}

\section*{De tangentibus Linearum curuarum.}

Ad superiorem methodum, inuentionem tangentium ad data puncta in lineis quibuscumq; curuis reducimus.

\note{Figure, dans la marge gauche sur toute la hauteur du texte : une parabole de sommet d, ouverte vers le bas, dont une branche passe par B et l'autre par n ; l'axe descend de d vers c ; deux ordonnées horizontales, O---I plus haut et B---C plus bas ; au-dessus de d, sur le prolongement de l'axe, le point E, d'où la droite EB descend, tangente en B, en passant par O.}

Sit data verbigratia parabole Bdn, cuius vertex d, diameter dc, \& punctum in eâ datum B ad quod ducenda est recta BE tangens parabolen Et in puncto E cum diametro concurrens Ergo sumendo quodlibet punctum O in recta BE, Et ab eo ducendo ordinatam OI, a puncto autem B ordinatam BC. maior erit proportio cd ad di quam quadrati BC ad quadratum $\overline{OI}$ quia punctum O est extra parabolen. Sed propter similitudinem triangulorum ut BC quadratũ ad $\overline{OI}$ quadratum ita CE quadratum ad $\overline{IE}$ quadratum. Maior igitur erit proportio $\overline{cd}$ ad di quam quadrati $\overline{CE}$ ad quadratum $\overline{IE}$. Cum autem punctum B detur, datur applicata BC. Ergo punctum $\overline{C}$. datur etiam $\overline{cd}$. Sit igitur cd æqualis * B data. Ponatur CE esse A. ponatur
\marginal{* Il a icy nommé B ce qu'il nomme d par apres}
\note{Le « O » de « quodlibet punctum O » est écrit au-dessus de la ligne. « ordinatam OI » : les deux premières lettres sont repassées d'une encre plus noire ; une tache mange la fin de la ligne (« Et ab eo »). La note marginale, en français, à gauche de la dernière ligne et appelée par l'astérisque devant « B data », parle de l'auteur à la troisième personne : c'est un lecteur ou le copiste, non l'auteur. La phrase continue en tête de la vue 30.}

\page{30}
\note{Verso du feuillet de la vue 29 ; aucun nombre en tête. La page continue la phrase de la vue 29 ; la figure de la vue 29 se voit par transparence dans la marge gauche.}

$\overline{CI}$ esse E. Ergo d ad d $-$ E. habebit maiorem proportionem quàm Aq. ad Aq $+$ Eq $-$ A in E$''$. Et ducendo inter se medias et extremas. d in Aq. $+$ d in Eq $-$ d in A in E$''$ maius erit quàm d in A quadr. $-$ Aq in E. Adæquentur igitur iuxta superiorem methodum. demptis itaq; communibus d in Eq $-$ d in A in E$''$ adæquabitur $-$ Aq in E. Aut (quod idem est. d in Eq $+$ Aq in E adæquabitur d in A in E$''$. omnia diuidantur per $\overline{E}$. ergo d in E $+$ Aq adæquabitur d in A$''$. Elidatur d in E. Ergo Aq æquabitur d in A$''$. Ideoq; A æquabitur d bis. Ergo CE probauimus duplam ipsius cd quod quidem ita se habet.
\note{« d in Eq $-$ d in A in E$''$ adæquabitur » : un petit trait vertical est posé sur le signe moins, comme sur le moins de la vue 28 ; on lit un moins, que la ligne suivante confirme en changeant les signes. Deux petites croix à l'encre sont dans la marge gauche, au droit de la quatrième et de la septième ligne ; leur sens n'est pas vu. La parenthèse ouverte avant « quod idem est » n'est pas refermée.}

Nec vnquam fallit methodus. Imo ad plerasq; quæstiones pulcherrimas potest extendi. Eius enim beneficio centra grauitatis in figuris, lineis curuis et rectis comprehensis Et in solidis inuenimus. Et multa alia de quibus fortasse alias si otium suppetat. De quadraturis spatiorum sub lineis curuis et rectis contentorum, imò et de proportione solidorum ab eis ortorum ad conos eiusdem basis et altitudinis fusè iam cum Dño de Roberual egimus
\note{Le texte s'arrête sur « egimus », sans point, au tiers inférieur de la page ; le reste est blanc. C'est la fin de la pièce ouverte à la vue 29. Dans la marge gauche, écrit de bas en haut, au crayon gris et d'une autre main : « Mr des Cartes \uncertain{f.} 347 », renvoi dont l'objet n'est pas dit sur le feuillet. Cachet rond de la Bibliothèque nationale au pied.}

\page{31}\folio{12r}
\note{Nouvelle pièce et nouvelle main (voir l'en-tête). En haut à droite : « 12 » au crayon rouge, au bout de la première ligne du titre ; au-dessous « 12 » au crayon gris. Au pied, au centre, « 43 » à l'encre (voir l'en-tête). Cachet rond de la Bibliothèque nationale dans la marge gauche.}

\section*{Appendix ad Isagogen topicam continens Soluõnem problematum Solidorum per locos.}
\note{« Soluõnem » : « Solutionem » abrégé, tilde sur la ligne.}

Patuit methodus quâ lineæ locales deteguntur. Inquirendum restat quarãone problematum Solidorum Solutio possit ex Supradictis elegantissimè deriuari. Hoc Vt fiat, coarctanda illa quantitatum ignotarum extra limites Suos euagandi licentia. Infinita enim Sunt puncta quibus quæstioni propositæ Satisfit in locis, cõmodissime igitur per duas æqualitates locales quæstio determinatur. \struck{\ill{}} secant Quippe se inuicem duæ lineæ locales positione datæ, et punctum in\struck{tersectioni}s \struck{\ill{}}Sectionis positione datum quæstionem ex Infinito ad terminos præscriptos adigit.
\note{« quarãone » : « qua ratione ». La fin de la phrase est corrigée à plusieurs endroits. En bout de ligne, après « determinatur. », un mot est barré (il commence peut-être par « Inter- ») et « secant » est écrit au-dessus. « inuicem » est ajouté au-dessus de la ligne, avec un signe d'insertion, entre « se » et « duæ ». En bout de ligne suivante, après « punctum in », le reste du mot (« tersectioni ») est barré et un « s » écrit au-dessus ; au début de la ligne d'après, un mot barré, illisible, précède « Sectionis ». La leçon voulue paraît « Secant quippe se inuicem … et punctum intersectionis positione datum ».}

Exemplis breuiter et dilucidè res explicatur

Proponatur A cubus $+$ B in A quadratum \struck{æquari}

æquari Z plano in B
\note{Les deux lignes sont disposées en tableau : « Proponatur » et « æquari » à gauche, les termes en retrait. Le premier « æquari », en bout de ligne, est barré et récrit à la ligne suivante.}

Commodè Vtraq; æqualitatis pars potest æquari Solido B in A in E. Vt per diuisionem istius Solidi Illinc per $\overline{A}$, hinc per B res deducatur ad locos, Cum igitur A cubus $+$ B in A quad. æquetur B in A in E, ergo Aq $+$ B in A æquabitur B in E.

Et erit Vt patet ex nrâ methodo extremitas ipsius E ad Parabolam positione datam.

Deinde cum Zpl. in B æquetur B in A in E. ergo Zpl. æquabitur A in E.

Et erit ex nrã methodo, extremitas ipsius $\overline{E}$ ad hyperbolam positione datam. Sed iam probauimus esse ad parabolam positione datam. Ergo datur positione, et est facilis ab analysi ad Synthesim regressus.

Nec dissimilis est methodus in omnibus æquationibus Cubicis. Constitutis enim ex Vna parte Solidis omnibus ab A adfectis ex alterâ Solido omnino dato Vel etiam cum Solidis ab A Vel Aq affectis poterit fingi æqualitas Superiori Similis

Proponatur exemplum in æquationibus quadrato quadratorum.

Aqq. $+$ B Sol. in A $+$ Zpl. in Aq æquetur D\struck{pl}ppl. ergo Aqq æquabitur Dppl. B Sol. in A Zq in Aq. æquentur hæc duo homogenea Zq in Eq.
\note{Entre « Dppl. », « B Sol. in A » et « Zq in Aq » aucun signe n'est écrit ; on laisse les blancs. Les Z de la page (« Z plano », « Zpl. », « Zq », « Z in E ») sont d'une encre plus noire et d'un tracé plus lourd que le reste, comme remplis après coup dans des blancs.}

Cum igitur Aqq \struck{æquabitur Dppl. \ill{}} æquetur Zq in Eq

Ergo per Subdiuisionem quadraticam

Aq. æquabitur Z in E. et erit extremitas $\overline{E}$ ad parabolam positione datam.
\note{Dans la marge droite, au droit de « æquetur Zq in Eq », un petit signe en forme de « u » ou de « n », dont le sens n'est pas vu. La page finit sur « positione datam. » ; la vue 32 continue le même développement.}

\page{32}
\note{Verso du feuillet de la vue 31 ; aucun nombre. Même main que la vue 31 ; la page continue son développement.}

deinde cum Dppl. B Solid. in A $-$ Zq. in Aq æquetur Zq in Eq omnibus per Zq diuisis

$\dfrac{\text{Dppl.}}{\text{Zq}} - \dfrac{\text{B Sol. in A}}{\text{Zq}} - \text{Aq}$ æquabitur Eq.

\struck{Eq}
\note{Entre « Dppl. » et « B Solid. » la première ligne n'a pas de signe. Dans la ligne suivante, « Dppl. » et « B Sol. in A » sont chacun écrits sur un trait, avec « Zq » dessous ; un trait court suit « Eq. ». Au-dessous, à gauche, un « Eq » isolé est barré.}

Et erit, ex nrã methodo extremitas $\overline{E}$ ad circulum positione datum, Sed est et ad parabolam positione datam. ergo datur.

Non dissimili methodo Soluentur quæstiones omnes quadrato quadraticæ. Expurgabuntur enim methodo Vietæ cap. i. de emend. ab affectione Sub Cubo; et quadrato quadrato ignoto ab Vna parte, reliquis homogeneis ab alterâ constitutis, per parabolam, circulum Vel hyperbolam Soluetur quæstio.

Proponatur ad exemplum inuentio duarum mediarum in continuâ proportione.

Sint duæ rectæ B \uncertain{maior} D minor inter quas duæ mediæ proportionales Sunt inueniendæ. fiet

A cubus æqualis Bq in D, posito \uncertain{nempe quod}

* major mediarum ponatur A
\note{Le premier des deux adjectifs après « B » et « D » est tracé presque comme le second (« minor ») ; la lecture « maior » est offerte d'après la suite (A, la plus grande des moyennes, avec A cube égal à Bq in D). Les derniers mots de la ligne sont rapides. L'astérisque devant « major », très appuyé, est dans la marge gauche.}

æquentur Singula homogenea B in A in E.

Illinc fiet Aq æquale B in E

Istinc A in E æquale B in D

Ideoq; quæstio per hyperbol\uncertain{es} et parabol\uncertain{es} intersectionem perficietur
\note{La fin de « hyperbol- » est sous un pâté d'encre ; celle de « parabol- » est surchargée.}

\note{Figure : une horizontale O---V---N ; en O, en V et en N s'élèvent trois perpendiculaires, terminées en R, M et Z ; la lettre V est un petit u. À droite, au-dessus, deux segments horizontaux marqués B et d, les droites données.}

Exponatur enim recta quæuis positione data OVN in quâ detur punctum O, Sint rectæ datæ B et D inter quas duæ mediæ proportionales inueniendæ. Ponatur recta $\overline{OV}$ æquari $\overline{A}$, Et recta $\overline{VM}$ ipsi $\overline{OV}$ ad rectos angulos, æquari $\overline{E}$. Ex priori æqualitate quâ Aq æquatur B in E, constat per punctum $\overline{O}$ tanquam Verticem describendam parabolam cuius rectum latus Sit B, diameter ipsi VM parallela, et applicatæ ipsi OV: transibit igitur hæc parabola per punctum m. Ex 2\textsuperscript{a}. æqualitate, quâ
\note{Plusieurs mots sont soulignés d'un trait fin (« quæuis », « Sitione » dans « positione », « media », « Verticem ») ; un crochet à l'encre encadre « Ex 2\textsuperscript{a}. » au bout de l'avant-dernière ligne. La phrase continue en tête de la vue 33.}

\page{33}\folio{13r}
\note{En haut à droite : « 13 » au crayon gris et « 13 » au crayon rouge, côte à côte. Au pied, au centre, « 44 » à l'encre (voir l'en-tête). Dans la marge gauche, au crayon gris, un signe « /n » au droit de la troisième ligne. La page continue la phrase de la vue 32 (« Ex 2\textsuperscript{a}. æqualitate, quâ »).}

B in d æquatur A in E. Sumatur punctum Vbilibet in rectâ OV ut n a quo excitetur perpendicularis VZ, et fiat rectan.\textsuperscript{lum} OnZ æquale rectangulo B in D. Excitetur etiam perpendicularis OR. Circa \struck{\ill{}} asymptotos RO, OV describenda hyperbola per punctum $\overline{Z}$ ex nrã methodo locali dabitur positione et transibit per punctum m. Sed parabol\uncertain{a} etiam quam Supra descripsimus datur positione, et per Idem punctum m transit, datur igitur punctum m positione, A quo si demittatur perpendicularis mV dabitur punctum V Et recta OV, major duarum continuè proportionalium quas quærimus
\note{« perpendicularis VZ » : on donne les lettres telles qu'elles sont tracées (le premier signe est le petit u, V) ; la figure de la vue 32 élève la perpendiculaire en n. Après « Circa », un mot court est barré. La fin de « parabol- » est surchargée. Mots soulignés d'un trait fin : « Vbilibet », « Excitetur », « Circa », « Supra », « positione », « Idem ».}

Inuenta igitur Sunt duæ mediæ per InterSectionem parabolæ et hyperbolæ

Si ad quadratoquadrata lubeat quæõnem extendere omnia ducantur in A.

Aqq. æquabitur Bq in D. in A.

Aequentur Singula homogenea iuxta Superiorem methodum Bq in Eq.

fient duæ æqualitates nempe Aq et B in E.

et D in A et Eq
\note{Un trait courbe au-dessus de la ligne relie « Aq et B in E » à « D in A et Eq », écrit en retrait sous la précédente.}

Quæ Singula dabunt parabolam positione datam. fiet igitur constructio mesolabij per Intersectionem duarum parabolarum hoc casu

Prior constructio et posterior Sunt apud Eutocium in Archimedẽ et huic methodo facillimè redduntur obnoxiæ.

Abeant igitur illæ parapleroses Vietææ quibus æquaõnes quadratoquadraticas reducit ad quadraticas per medium Cubicarũ abs radice planâ. pari enim elegantia facilitate et breuitate Soluuntur, Vt iam patuit, perinde quadratoquadraticæ ac Cubicæ quæõnes, nec possunt, opinor, elegantius.
\note{Après « Abeant igitur », un blanc d'environ un mot, marqué d'un trait fin. « æquaõnes », « quæõnes » : « æquationes », « quæstiones » abrégés.}

Vt pateat elegantia hujus methodi. En constructionem oĩum problematum cubicorum et quadratoquadraticorum per parabolam et circulum

Ponatur Aqq. $+$ Z Sol. in A æquari Dppl. ergo Aqq æquabitur Z Sol. in A $+$ d ppl. Fingatur quadratum abs Aq $-$ Bq aut alio quouis quadrato dato. fiet quadratum, Aqq $+$ Bqq $-$ Bq in Aq bis. Addantur ad Supplementum Singulis æqualitatis partibus Bqq $-$ Bq in Aq bis. fiet Aqq $+$ Bqq, Bq in Aq bis
\note{Le copiste termine souvent ses lignes par un tiret de remplissage. Trois tirets de bout de ligne tombent ici entre deux termes : après « ergo Aqq », après « abs Aq » et après « Aqq $+$ Bqq ». On lit les deux derniers comme des moins, que la suite demande ; le premier est laissé pour un tiret de remplissage, et « æquabitur Z Sol. in A $+$ d ppl. » est donné tel qu'il est écrit. Entre « Bqq, » et « Bq in Aq bis » à la dernière ligne, aucun signe. Les Z sont encore d'une encre plus noire. La phrase continue en tête de la vue 34.}

\page{34}
\note{Verso du feuillet de la vue 33 ; aucun nombre. La page continue la phrase de la vue 33.}

fiet Aqq $+$ Bqq $-$ Bq in Aq bis æquale. Bqq $-$ Bq in Aq $-$ Z sol. in A $+$ Dppl. Sit Bq bis æquale Nq Et Singulis homogeneis Siue partibus æqualitatis æquetur Nq in Eq. fiet Illinc per Subdiuisionem quadraticam Aq $-$ Bq æquale N in E, Ideoq; punctum extremum E erit ad parabolam ex nrã methodo, Istinc fiet
$\dfrac{\text{Bqq}}{\text{Nq}} - \text{Aq} - \dfrac{\text{Z sol. in A}}{\text{Nq}} + \dfrac{\text{Dppl.}}{\text{Nq}}$ æquale Eq.
\note{La lettre qu'on donne N est écrite comme un n minuscule, qu'on peut confondre avec le petit u (V) du même copiste ; en bas de la page elle est nettement un n. Les fractions sont écrites comme ici, numérateur sur un trait, « nq » dessous ; un petit trait oblique suit « Eq. ».}

Ideoq; ex nrã methodo punctum extremum E erit ad circulum descriptione igitur parabolæ et circuli Soluitur quæõ.

Hæc methodus facillimè ad omnes Casus tam Cubicos quam quadroquadraticos extenditur. curandum enim tantum Vt ex Vna parte Sit Aqq, Ex altera quælibet homogenea modo non afficiantur ab A cubo. at per expurgaõnem Vietæam omnes æquaõnes quadratoquadraticæ ab affectione Sub Cubo liberantur, Ergo Eadem in omnibus methodus. Cum autem æquationes cubicæ liberentur ab adfectione Sub quadrato per methodum Vietæam, homogeneis oĩbus in A ductis, fiet æquatio quadratoquadratica, cuius nullum ex homogeneis afficietur Sub Cubo, ideoq; Soluetur per Superiorem methodum.

Id Solum in 2\textsuperscript{a}. æqualitate curandum est, Vt Aq ex Vna parte ex altera Eq Sub contraria affectionis notâ reperiantur, quod est Semper facillimum
\note{Dans la marge gauche, au droit de ce paragraphe, un long trait ondulé vertical à l'encre, et un tiret devant la ligne suivante ; ni l'un ni l'autre n'est du texte.}

Sit Enim in alio casu, Vt omnia percurramus

Aqq æquale Zpl. in Aq $-$ Z Sol. in d

fingatur quoduis quadratum abs Aq $-$ quouis quadrato dato Vt Bq fiet Aqq $+$ Bqq $-$ Bq in Aq bis. Adjiciatur Vtrique æqualitatis parti ad Supplementum Bqq $-$ Bq in Aq bis.

fiet Aqq $+$ Bqq $-$ Bq in Aq bis æquale

Bqq $-$ Bq in Aq bis $+$ Zpl. in Aq $-$ Zsol. in d.

Vt Igitur commoda fiat diuisio in 2\textsuperscript{a}. æqualitate Sumenda differentia inter Bq bis \& Zpl. quæ sit Verbigratiâ Nq et Vtraq; æqualitatis pars æquanda Nq in Eq

Vt Illinc fiat Aq $-$ Bq æquale N in E.

Istinc
$\dfrac{\text{Bqq}}{\text{Nq}} - \text{Aq} - \dfrac{\text{Z Sol.}}{\text{Nq}}$ in d. æquale Eq
\note{Dans cette ligne, « in d. » est écrit hors du trait de fraction, à la suite de « Z Sol. » ; un pâté couvre la fin de « Eq ».}

Aduertendum deinde Bq bis debere præstare Z plano, alioquin
\note{Deux points sont tracés au-dessus du Z de « Z plano ». La phrase continue en tête de la vue 35.}

\page{35}\folio{14r}
\note{En haut à droite : un nombre au crayon gris, de lecture douteuse, « 16 » ou « 13 » (le second chiffre est une boucle ouverte), et « 14 » au crayon rouge. Au pied, au centre, « 45 » à l'encre (voir l'en-tête). La page continue la phrase de la vue 34 (« alioquin »).}

Aq non afficeretur Signo defectus, Et pro circulo inueniremus hyperbolam, cui promptum remedium, Bq enim ad libitum Sumimus, Ideoq; ipsius duplum majus Zpl. nullius est negotij Sumere. Constat autem ex methodo locali circulum creari Semper ex æqualitate in cuius parte altera quadratum Vnum ignotum afficitur Signo \struck{plus} $+$, in alterâ, aliud quadratum ignotum, Signo \struck{minus} $-$.
\note{« plus » et « minus » sont barrés, et les signes « $+$ » et « $-$ » écrits au-dessus.}

Si Sumas ad hoc exemplum inuentionem duarum mediarum erit AC. æqualis. Bq in d.

Et Aqq æquale Bq in d in A

Adjiciatur Vtrinq; Bqq $-$ Bq in Aq$''$

Aqq $+$ Bqq $-$ Bq in Aq$''$

æquabitur Bqq $+$ Bq in d in A $-$ Bq in Aq$''$

Sit Bq$''$ æquale Nq.

Et Singulæ æqualitatis partes æquentur Nq in Eq

fiet Illinc Aq $-$ Bq æquale N in E.

Ideoq; Extremum E erit ad parabolam

Istinc fiet Bq $\frac{1}{2}$ $+$ d $\frac{1}{2}$ in A $-$ Aq æquale Eq. ideoq; extremum E erit ad circulum
\note{« AC. » : « A cubus » abrégé. Les deux traits après « Aq » et « Bq » valent « bis », comme dans les pièces précédentes. « Bq » et « d » sont suivis chacun d'une petite fraction « $\frac{1}{2}$ » écrite sous la ligne ; « d » est souligné. Un long trait courbe sous « extremum E erit ad circulum » sépare la conclusion.}

Qui hæc aduerterit frustrà quæõnem mesolabij, trisectionis angularis, et Similes tentabit deducere ex planis, hoc est per rectas \& circulos expedire.
\note{Fin de la pièce ouverte à la vue 31 ; le reste de la page est blanc. Cachet rond de la Bibliothèque nationale sous la dernière ligne.}

\page{37}\folio{15r}
\note{Nouvelle pièce, sur un nouveau feuillet ; la main paraît être celle des vues 31–35. En haut à droite : « 14 » au crayon gris et, au-dessous, « 15 » au crayon rouge. En haut à gauche, à l'encre, une lettre isolée « E » (ou un epsilon), qui n'est pas du texte. Cachet rond de la Bibliothèque nationale dans la marge gauche.}

M\textsuperscript{r}. fermat
\note{En tête, au centre, d'une encre plus pâle et d'un trait plus fin que le texte, peut-être d'une autre main : « Mr. fermat », le « r » en exposant et souligné. C'est la première fois dans ce lot qu'un feuillet porte le nom de Fermat. C'est une inscription en tête de la pièce, qui la rapporte à lui ; elle ne dit pas que la main est la sienne, et on ne l'affirme pas.}

Exponatur tres rectæ positione datæ triangulum constituentes. AM, MB, BA. Et Sit quoduis punctum O à quo ad rectas datas ducantur rectæ OE, OI, OD, in angulis OEM, OIM, ODB, datis, Sit autem ratio rectanguli EOD ad quadratum OI data. Aio punctum O esse ad Vnam ex Coni Sectionibus. Diuidatur MB bifariam in q, et iuncta Aq, ducantur per punctum O rectæ foc, on ipsis MB, MA parallelæ Tria Triangula oef, odc, oin Sunt Specie data. nam ex hypothesi dantur Anguli oef, odc, oin Dantur etiam efo quia propter parallelas dato AMB est æqualis, Datur et ocd quia æqualis dato mbA. Denique datur oni, cum detur onB ipsi AMB propter parallelas æquales. Datur Igitur ratio oE ad of, datur Item ratio od ad oc, Ergo ratio rectanguli eod ad rectang.\textsuperscript{um} foc datur. datur autem ex hypothesi ratio rectanguli eod ad quadratum oi, ergo ratio rectanguli foe ad quadratum oi datur, Datur autem ratio quadrati oi ad quadratum on, propter datum \struck{\ill{}} Specie triangulum oin, Ergo ratio rectanguli foc ad quadratum on Siue fm ipsi æquale datur. Si \uncertain{fertur} Aq in V, ita Vt ducatur parallelâ mb quadratum VR ad quadratum rm Sit in ratione datâ rectanguli foc ad quadratum fm (hoc autem est facile, cum angulus mrV detur) et per punctum V describatur circa diametrum Aq Coni Sectio, quam rectæ ma, ab in punctis mb contingant (Id autem est facillimum, et ex puncti V Situ erit aut parabole, aut hyperbole, aut Ellipsis. Superflua præsertim tibi non addimus) aio Coni Sectionem Sic descriptam per punctum O transire: nam transeat ex aliâ parte per punctum p, tanget recta Vr Sectionem cum Sit parallela ordinatæ mb. Cum igitur Sectio transeat per punctum O, erit rectangulum pfo ad quadratum fm. Vt quadratum Vr ad quadratum rm \struck{Ideoque est rectangulum} ex sexta decima propoõne 3\textsuperscript{i}. Apoll. Vt autem quadratum Vr ad quadratum rm ita est rectang.\textsuperscript{um}
\note{Les lettres de la figure sont écrites tantôt en capitales, tantôt en minuscules (« M » et « m », « E » et « e ») ; on les donne comme elles sont écrites ; la lettre donnée V est le petit u du copiste. Aucune figure n'est tracée sur la page. Beaucoup de mots et de lettres sont soulignés d'un trait fin, comme aux vues 32–33 ; on ne le rend pas. « foe » (ligne 12) est écrit ainsi, pour le « foc » des lignes voisines. « propter datum » est suivi d'un mot court barré. « Si \uncertain{fertur} Aq in V » : le premier mot se lit « fertur » (ou « sertur ») ; le sens attend « secetur » ou « feratur ». Après « et ex », un blanc d'un mot est laissé dans la ligne. À l'avant-dernière ligne, « Ideoque est rectangulum » est barré, et « decima » est écrit au-dessus de « sexta » : la leçon finale paraît « ex sexta decima propoõne 3\textsuperscript{i}. Apoll. », la proposition 16 du livre III d'Apollonius. « propoõne » : « propositione » abrégé. La phrase continue au verso, vue 38.}

\page{38}
\note{Verso du feuillet de la vue 37 ; aucun nombre. La page continue la phrase de la vue 37 (« ita est rectang.\textsuperscript{um} »).}

foc ad quadratum fm ex Constructione. rectang.\textsuperscript{um} igitur pfo rectangulo foc æquale erit, Ideoque recta fo, rectæ pc. Quod quidem ita se habet; nam cùm Aq diuidat bifariam mb, erit recta fx rectæ xc æqualis, propter Sectionem Vero recta ox est æqualis xp, reliqua igitur fo rectæ pc æquatur. nec est difficilis ab analysi ad Synthesim per demonstraõnem ducentem ad impossibile regressus.
\note{« demonstraõnem » : « demonstrationem » abrégé. Le texte s'arrête sur « regressus. », au tiers supérieur de la page ; le reste est blanc, la vue 37 se voyant par transparence. Cachet rond de la Bibliothèque nationale sous la dernière ligne.}

Pour Mons.\textsuperscript{r} \uncertain{Carcaui}. rue michel le Conte au milieu
\note{Adresse, écrite en travers, le long du bord droit de la moitié inférieure de la page, perpendiculairement au texte : le feuillet a été plié et porté comme une lettre. Le nom commence par une grande capitale bouclée, lue « C », suivie de « arcaui » ; la lecture « Carcaui » est offerte avec doute, sans rien en tirer de plus. La main de l'adresse paraît plus libre que celle du texte ; on ne décide pas si c'est la même. L'adresse ne nomme pas l'expéditeur.}

\page{40}
\note{Verso du feuillet blanc de la vue 39 ; aucun nombre, aucun texte. La page ne porte qu'une figure, à l'encre, dans son tiers supérieur, dont la lettre est celle de la pièce des vues 37–38 : c'en est sans doute la figure, que la page de texte n'a pas. Un triangle a, m, b (a en haut, m à gauche et b à droite sur la base) ; q au milieu de mb, et la droite aq ; un point o dans le triangle, par lequel passent la parallèle f---o---x---p---c à mb (f sur am, c sur ab) et la parallèle on à ma (n sur mb) ; de o partent oe vers am, od vers ab et oi vers mb (i entre m et n) ; u sur aq, et la parallèle r---u à mb, r sur am ; une courbe passe par u, o et p et descend jusqu'en b, et se prolonge en pointillé vers m. Deux petits crochets, en marge à gauche et à droite, au niveau de la base, ne sont pas des lettres de la figure.}

\end{document}
